\section{Тест производительности}
{\itshape Проводится сравнение скорости поиска образцов с помощью суффиксного массива и алгоритма Кнута-Морриса-Пратта (КМП). Время построения суффиксного массива при сравнении не учитывается, поскольку предполагается, что массив строится однократно для текста, а поиск выполняется многократно.}

Тест производительности представляет из себя следующее: генерируется текст заданной длины, состоящий из повторяющейся буквы \enquote{a}, а также 200 образцов со случайными длинами от 2 до 100 букв (буквы от \enquote{a} до \enquote{z}). Фиксируется время поиска всех образцов с помощью реализации на основе суффиксного массива и с помощью алгоритма КМП. Для получения статистически значимых результатов тесты проводятся для текстов размером 10\,000, 100\,000 и 1\,000\,000 символов. Время построения суффиксного массива замеряется отдельно и не включается в сравнение.

Для автоматизации тестирования использовался скрипт \texttt{benchmark.sh}:

\begin{alltt}
#!/bin/bash

echo "=== Building ==="
make bench
echo ""

run_test() \{
    local size=\$1
    local size_label=\$2

    echo "=== Testing with text size: \$\{size_label\} ==="

    ./gen.out \$size > test_bench.txt

    ./kmp.out < test_bench.txt > res_kmp.txt
    echo "KMP:"
    cat res_kmp.txt | grep "time"

    ./main.out < test_bench.txt > res_main.txt
    echo "Suffix array:"
    cat res_main.txt | grep "time"

    echo ""
\}

run_test 10000 "10,000"
run_test 100000 "100,000"
run_test 1000000 "1,000,000"

echo "=== All tests completed ==="
\end{alltt}

\noindent Результаты выполнения скрипта:

\begin{alltt}
root@6ec7ac378e08:/workspaces/discran-labs/laba-5# ./benchmark.sh
=== Building ===
g++ -O2 main.cpp -o main.out
g++ -O2 gen.cpp -o gen.out
g++ -O2 kmp.cpp -o kmp.out

=== Testing with text size: 10,000 ===
KMP:
KMP search time: 0.001486 sec
Suffix array:
Suffix array build time: 0.006544 sec
Suffix array search time: 0.000015 sec

=== Testing with text size: 100,000 ===
KMP:
KMP search time: 0.010319 sec
Suffix array:
Suffix array build time: 0.824346 sec
Suffix array search time: 0.000017 sec

=== Testing with text size: 1,000,000 ===
KMP:
KMP search time: 0.084405 sec
Suffix array:
Suffix array build time: 121.391583 sec
Suffix array search time: 0.000021 sec

=== All tests completed ===
\end{alltt}

\noindent Результаты тестов сведены в таблицу:

\begin{center}
\begin{tabular}{|c|c|c|c|}
\hline
\textbf{Размер текста} & \textbf{КМП (поиск)} & \textbf{Суф. массив (построение)} & \textbf{Суф. массив (поиск)} \\
\hline
10\,000 & 0.001486 с & 0.006544 с & 0.000015 с \\
\hline
100\,000 & 0.010319 с & 0.824346 с & 0.000017 с \\
\hline
1\,000\,000 & 0.084405 с & 121.391583 с & 0.000021 с \\
\hline
\end{tabular}
\end{center}

Как видно из таблицы, время поиска с помощью суффиксного массива практически не зависит от длины текста и остаётся на уровне сотых долей миллисекунды для всех размеров. В то же время время поиска КМП растёт линейно с ростом текста: от 0.0015 с для 10\,000 символов до 0.0844 с для 1\,000\,000 символов. Таким образом, суффиксный массив обеспечивает выигрыш в скорости поиска примерно в 100--4000 раз для данного набора тестов.

Однако следует отметить, что построение суффиксного массива в данной реализации имеет квадратичную сложность из-за использования стандартной сортировки с полным сравнением подстрок. Это приводит к резкому росту времени построения: 0.007 с для 10\,000 символов, 0.824 с для 100\,000 символов и 121.4 секунды для 1\,000\,000 символов. Поэтому применение такого подхода оправдано только в сценариях, где текст задан однократно и не изменяется, а количество запросов на поиск велико. В противном случае алгоритм КМП, не требующий предобработки текста (только препроцессинг образца за линейное время), оказывается более практичным выбором.

\pagebreak